\chapter{Communication Complexity --- Public Coin Basic}
%%% generated by the blueprint pipeline from TCSlib.CommunicationComplexity.NewmanTheorem.PublicCoinBasic
\section{Overview}

This module introduces public-coin communication protocols, where both Alice and Bob
have access to a shared source of randomness $\Omega$.  It defines the core type
\texttt{CommunicationComplexity.PublicCoin.Protocol}, the execution function
\texttt{CommunicationComplexity.PublicCoin.Protocol.rrun}, and two notions of approximate
correctness: $\varepsilon$-satisfying a predicate and $\varepsilon$-computing a function.

%%%%%%%%%%%%%%%%%%%%%%%%%%%%%%%%%%%%%%%%%%%%%%%%%%%%%%%%%%%%%%%%%%%%%%%%%%%%%%%
\section{Declarations}

\begin{definition}[Public-coin protocol]
\lean{CommunicationComplexity.PublicCoin.Protocol}
\statementsource{roughgarden-cc-bootcamp}{pdf-d54d402b16e2-p011-b003}
\statementsource{raoyehudayoff-ch03-randomized-protocols}{
  pdf-3eebfb23ab53-p003-b004,
  pdf-3eebfb23ab53-p004-b001
}
\label{CommunicationComplexity.PublicCoin.Protocol}
\leanok
\uses{CommunicationComplexity.Deterministic.Protocol}
A public-coin protocol over shared randomness $\Omega$, with Alice's private input
in $X$ and Bob's in $Y$ producing an output in $\alpha$, is a deterministic protocol
whose combined input types are $\Omega \times X$ (for Alice) and $\Omega \times Y$
(for Bob).  In other words, both players observe the same random string $\omega \in \Omega$
alongside their private inputs.
\end{definition}

\begin{definition}[Output node]
\lean{CommunicationComplexity.PublicCoin.Protocol.output}
\label{CommunicationComplexity.PublicCoin.Protocol.output}
\leanok
\uses{CommunicationComplexity.PublicCoin.Protocol}
The leaf node of a public-coin protocol that returns the constant value $a \in \alpha$,
constructed by lifting the corresponding deterministic output node.
\end{definition}

\begin{definition}[Alice's move]
\lean{CommunicationComplexity.PublicCoin.Protocol.alice}
\label{CommunicationComplexity.PublicCoin.Protocol.alice}
\leanok
\uses{CommunicationComplexity.PublicCoin.Protocol}
Constructs a public-coin protocol node at which Alice sends a single bit: given a
function $f : X \to \Omega \to \mathrm{Bool}$, Alice transmits $f\,x\,\omega$ and
the protocol continues according to $P : \mathrm{Bool} \to \mathrm{Protocol}$.
\end{definition}

\begin{definition}[Bob's move]
\lean{CommunicationComplexity.PublicCoin.Protocol.bob}
\label{CommunicationComplexity.PublicCoin.Protocol.bob}
\leanok
\uses{CommunicationComplexity.PublicCoin.Protocol}
Constructs a public-coin protocol node at which Bob sends a single bit: given a
function $f : Y \to \Omega \to \mathrm{Bool}$, Bob transmits $f\,y\,\omega$ and
the protocol continues according to $P : \mathrm{Bool} \to \mathrm{Protocol}$.
\end{definition}

\begin{definition}[Randomized execution]
\lean{CommunicationComplexity.PublicCoin.Protocol.rrun}
\statementsource{raoyehudayoff-ch03-randomized-protocols}{
  pdf-3eebfb23ab53-p003-b004,
  pdf-3eebfb23ab53-p004-b001
}
\label{CommunicationComplexity.PublicCoin.Protocol.rrun}
\leanok
\uses{CommunicationComplexity.Deterministic.Protocol.run, CommunicationComplexity.PublicCoin.Protocol}
Given a public-coin protocol $p$, private inputs $x \in X$, $y \in Y$, and a shared
random string $\omega \in \Omega$, $\texttt{rrun}\,p\,x\,y\,\omega$ executes the
underlying deterministic protocol on the paired inputs $(\omega, x)$ and $(\omega, y)$
and returns the resulting output.
\end{definition}

\begin{theorem}[Randomized execution as deterministic execution on paired inputs]
\lean{CommunicationComplexity.PublicCoin.Protocol.rrun_eq}
\label{CommunicationComplexity.PublicCoin.Protocol.rrun_eq}
\leanok
\uses{CommunicationComplexity.Deterministic.Protocol.run, CommunicationComplexity.PublicCoin.Protocol, CommunicationComplexity.PublicCoin.Protocol.rrun}
\difficulty{1}
Let $p$ be a public-coin protocol with shared randomness in $\Omega$, Alice's private
input in $X$, and Bob's in $Y$, producing an output in $\alpha$. Then for every $x \in
X$, $y \in Y$, and $\omega \in \Omega$, the randomized execution of $p$ on $x$, $y$ with
random string $\omega$ equals the execution of the underlying deterministic protocol on
the paired inputs $(\omega, x)$ and $(\omega, y)$.
\end{theorem}

\begin{definition}[$\varepsilon$-satisfies a predicate]
\lean{CommunicationComplexity.PublicCoin.Protocol.ApproxSatisfies}
\statementsource{raoyehudayoff-ch03-randomized-protocols}{pdf-3eebfb23ab53-p004-b001}
\label{CommunicationComplexity.PublicCoin.Protocol.ApproxSatisfies}
\leanok
\uses{CommunicationComplexity.PublicCoin.Protocol, CommunicationComplexity.PublicCoin.Protocol.rrun}
A public-coin protocol $p$ \emph{$\varepsilon$-satisfies} a predicate
$Q : X \to Y \to \alpha \to \mathrm{Prop}$ if for every input pair $(x, y)$,
\[
  \mathrm{vol}\bigl\{\omega \in \Omega \mid \neg\, Q\,x\,y\,(p.\texttt{rrun}\,x\,y\,\omega)\bigr\} \;\le\; \varepsilon.
\]
\end{definition}

\begin{definition}[$\varepsilon$-computes a function]
\lean{CommunicationComplexity.PublicCoin.Protocol.ApproxComputes}
\statementsource{roughgarden-cc-bootcamp}{pdf-d54d402b16e2-p011-b003}
\statementsource{raoyehudayoff-ch03-randomized-protocols}{pdf-3eebfb23ab53-p004-b001}
\label{CommunicationComplexity.PublicCoin.Protocol.ApproxComputes}
\leanok
\uses{CommunicationComplexity.PublicCoin.Protocol, CommunicationComplexity.PublicCoin.Protocol.rrun}
A public-coin protocol $p$ \emph{$\varepsilon$-computes} a function
$f : X \to Y \to \alpha$ if for every input pair $(x, y)$,
\[
  \mathrm{vol}\bigl\{\omega \in \Omega \mid p.\texttt{rrun}\,x\,y\,\omega \ne f\,x\,y\bigr\} \;\le\; \varepsilon.
\]
Here the volume is taken with respect to the measure on $\Omega$ given by the
\texttt{MeasureSpace} instance.
\end{definition}

\begin{theorem}[Approximate computation as satisfying pointwise equality]
\lean{CommunicationComplexity.PublicCoin.Protocol.ApproxComputes_eq_ApproxSatisfies}
\label{CommunicationComplexity.PublicCoin.Protocol.ApproxComputes_eq_ApproxSatisfies}
\leanok
\uses{CommunicationComplexity.PublicCoin.Protocol, CommunicationComplexity.PublicCoin.Protocol.ApproxComputes, CommunicationComplexity.PublicCoin.Protocol.ApproxSatisfies}
\difficulty{1}
Let $p$ be a public-coin protocol over a shared-randomness space $\Omega$ carrying a
measure, with Alice's inputs in $X$, Bob's in $Y$, and outputs in $\alpha$, let $f : X
\to Y \to \alpha$ be a function, and let $\varepsilon \in \bbr$. Then the assertion that
$p$ $\varepsilon$-computes $f$ coincides with the assertion that $p$
$\varepsilon$-satisfies the predicate $Q(x, y, a)$ given by $a = f(x, y)$; that is, the
two propositions are equal.
\end{theorem}
